% Chapter 26: Temperature Is Not Heat (complete chapter)
\chapter{Temperature Is Not Heat}

\step{A Counterintuitive Opening}

\begin{mainline}
Enter an unheated classroom in winter and touch an iron desk leg, then its wooden top. The iron feels icy, the wood mild, so different that you might blurt out, ``Iron is much colder than wood.'' Yet hold a thermometer against each and let its reading settle: their temperatures are identical. Both have spent the whole night in the same room. How could they differ?

Either your hand or the thermometer must be lying. Or, more interestingly, neither lies: they measure different things.

Consider another small puzzle. You test bathwater with a hand and find it just right. But after washing your hands in cold water, that same bath feels hot; after washing them in hot water, it feels cool. The water is unchanged; your hand has changed. When we say we feel hot or cold, what are we reporting?

% BEGIN TRANSLATOR SAFETY NOTE
\textbf{Translator safety note:} Do not let sparks land on skin or handle them to test this claim. Sparklers can cause serious burns, and welding requires protection from hot particles and radiation. See \href{https://www.cpsc.gov/node/65701}{CPSC sparkler guidance} and \href{https://www.osha.gov/etools/shipyard/general-requirements/ppe/ppe-selection/hot-work}{OSHA welding protection guidance}.
% END TRANSLATOR SAFETY NOTE

There are more extreme examples. Holiday sparklers and welding sparks reach over a thousand degrees Celsius, yet landing on your hand they often give only a mild sting or can be flicked away. A kettle of freshly boiled water is only a hundred degrees, but nobody would plunge a hand into it. Something at a thousand degrees can do less harm than something at a hundred. Clearly a distinction separates high temperature from harmful heating.

These examples point to one conclusion: our everyday word hot conflates three things. This chapter separates temperature, heat, and internal energy. Afterward you will see that saying a cup of water contains lots of heat is physically nonsensical, not merely imprecise: the sentence has a category error.
\end{mainline}

\step{Make a Guess First}

As usual, place your bets before the reveal. Write answers to these three questions.

First. Iron and wood have spent a night in the same room. Which has the lower temperature? (a) Iron, because it feels cooler. (b) Wood, because it cannot hold heat. (c) Equally low, or rather equally high: the same temperature.

Second. Which contains more heat, a drop of freshly boiled water or a large kettle of water at sixty degrees Celsius? (a) The boiling drop, because its temperature is higher. (b) The warm kettle, because there is more water. (c) The question itself is wrong.

Third. A cup contains water and unmelted floating ice. It sits on a table in a room that is not especially hot. After half an hour, half the ice has melted. What is the water temperature? (a) Higher, because the room keeps heating it. (b) Still zero degrees Celsius. (c) Slightly higher, by an incalculable amount.

Write your reasons too, and check them at the end. We are repairing the intuition that temperature, heat, and thermal sensation are one thing.

\step{Hands-on Experiment}

\begin{experiment}[Three Cups of Water and a Lying Hand]
% BEGIN TRANSLATOR SAFETY NOTE
\textbf{Translator safety note:} Do not follow the proposed sixty-degree-Celsius hand immersion or hold a hand in hot water for a minute. Water at that temperature can cause severe burns within seconds. Use a thermometer without skin immersion instead. See \href{https://www.cpsc.gov/s3fs-public/5098.pdf}{CPSC tap-water scald guidance}.
% END TRANSLATOR SAFETY NOTE

\textbf{Materials}: three identical cups; hot water around sixty degrees Celsius, with care against scalding; cold water with a few ice cubes; room-temperature water; a thermometer, either food or clinical; a metal spoon, wooden chopstick, and plastic ruler.

\textbf{Step one: three cups.} Put hot water in the left cup, cold water in the right, and warm water mixed half-and-half in the middle. Put your left hand in the hot water and your right in the cold, silently counting to sixty. Then put both hands simultaneously into the middle cup. Attend to their reports: the left says cool, the right says hot. Two hands in one cup give opposite answers.

\textbf{Step two: three materials.} Leave spoon, chopstick, and ruler side by side for half an hour to equalize temperatures. Close your eyes and briefly touch each with the back of your hand. Most people rank their coolness as metal first, plastic next, wood last.

\textbf{Step three: let the thermometer judge.} Hold it firmly against each surface for a minute or two. The readings are almost identical, all at room temperature.

\textbf{Step four: make a prediction.} Measure the hot water, perhaps fifty-five degrees, and cold water, perhaps five degrees. Pour half a cup of each, predict the mixed temperature mentally, then mix and measure. Most predict thirty degrees, and the result comes very close if you work quickly with cups that are not too thick. Record prediction and measurement. The mathematics will tell you when this equal-parts averaging works and when it fails.

\textbf{Record}: write both contradictions, opposite hand judgments of one cup and disagreement between hand and thermometer on one metal object. No explanation yet, but remember: the hand is easily recalibrated in step one and gives different readings for equal-temperature objects in step two. It probably does not measure temperature.
\end{experiment}

Add some free everyday observations to your question list. Rice microwaved for two minutes burns your mouth, while its ceramic bowl is merely warm. Why the difference in one oven? When measuring a fever, doctors insist on keeping the thermometer in place for five minutes. Why no instant reading on skin contact? Breathe on a cold window in winter and it fogs immediately; breathe on a mirror and the fog disappears in seconds. By chapter's end, these thermal stories should be clear.

\step{The Old Model Runs into Trouble}

State the old model fairly. History's most respectable answer to what heat is was \textbf{caloric theory}: heat is an invisible, massless fluid called caloric. Hot objects are rich in it, cold ones poor. Contact lets it flow from rich to poor until their levels equalize. Frictional heating squeezes hidden caloric out of objects.

This model was far more successful than its reputation suggests. It neatly explained thermal equalization after contact. It supported calorimetry, measurements of materials' capacities to hold caloric, and mixed-temperature calculations matching experiment exactly. Many thermal-engineering formulas still bear its outline today. It was an exceptionally useful wrong theory.

But fatal trouble arrived. In 1798, Count Rumford supervised cannon boring with a blunt drill at a Munich arsenal. Violent friction heated drill and barrel enough to warm nearby cold water. Caloric theory said friction merely squeezed out a finite stock of caloric. Yet the drill was so blunt it cut almost no chips, hence no new surfaces from which caloric could be squeezed, while heat kept pouring out as long as horses supplied motion. Heat was not released inventory but something that could be \emph{continually produced}. A supposedly conserved fluid appearing from nowhere was no fluid. Around the same time, Davy performed an even sharper experiment: rubbing two pieces of ice together in an evacuated vessel melted them. With no surrounding heat source or incoming caloric, the heat had to come from friction. A second nail entered caloric theory's coffin.

A second problem comes from your hands. One warm cup feels cool to the left hand and hot to the right. If the hand measures temperature, it cannot even produce equal outputs for equal inputs; its reading depends on its recent history. Either sensation is unreliable, or its actual job is not reporting temperature.

A third problem hides in language. We say hot water contains heat, yet rubbing hands warms them without anyone transferring heat: \emph{work} creates the warming. Press them against a cold wall and heat flows away again. Something contained in an object should not appear and vanish like that. The common diagnosis: heat is not a substance or inventory. We need a scale independent of hands, then a new grammatical role for heat.

\step{Inventing New Concepts}

\begin{mainline}
\textbf{New concept one: thermal equilibrium.} Put two objects of different hotness in contact and watch macroscopic properties: volume, resistance, color, anything temperature-sensitive. Both initially change, but eventually stop, reaching a macroscopically unchanging state called \emph{thermal equilibrium}. This is experimental fact: regardless of material, shape, or contact arrangement, sufficient waiting brings equilibrium.

\textbf{New concept two: the zeroth law of thermodynamics.} Let object A equilibrate with thermometer C, then B with C, without A and B ever touching. Now bring A and B together: nothing happens, because they are already in equilibrium. As a law: if A is in thermal equilibrium with C, and B with C, then A is in equilibrium with B. It sounds trivial, like the physical version of things equal to the same thing being equal. Yet it is an \emph{experimental law}, summarizing countless contact experiments, not a logical command nature must obey. Its role is so fundamental that, after the first and second laws had already been named, it had to be inserted as law zero. Its real purpose is authorization: thermometer C can \emph{represent} all objects equilibrated with it in comparisons. Temperature measurements can now communicate without A and B meeting.
\end{mainline}

\begin{center}
\begin{tikzpicture}[scale=1.0]
  % A, B, C and the zeroth law
  \draw[thick,fill=red!12] (0,0) circle (0.7);
  \node at (0,0) {A};
  \draw[thick,fill=blue!10] (4.4,0) circle (0.7);
  \node at (4.4,0) {B};
  \draw[thick,fill=yellow!25] (2.2,2.2) circle (0.7);
  \node at (2.2,2.2) {C (thermometer)};
  \draw[{Stealth}-{Stealth},thick] (0.55,0.62) -- (1.7,1.68);
  \node[above left] at (1.0,1.35) {\small Equilibrium};
  \draw[{Stealth}-{Stealth},thick] (2.7,1.68) -- (3.85,0.62);
  \node[above right] at (3.4,1.35) {\small Equilibrium};
  \draw[{Stealth}-{Stealth},thick,dashed] (0.7,0) -- (3.7,0);
  \node[below] at (2.2,-0.75) {\small Then A and B must be in equilibrium};
\end{tikzpicture}
\end{center}

\begin{mainline}
\textbf{New concept three: temperature.} The zeroth law tells us equilibrated objects share something. Name it \emph{temperature}. Notice the order: temperature does not precede equilibrium. The operational fact of equilibrium comes first; temperature is the shared label of an equilibrium group. Objects in one group have the same temperature, those in different groups different temperatures. Experiments further show that contact cools the initially hotter object and warms the cooler one, allowing labels to be ordered and scaled. A temperature scale assigns numbers: Celsius sets ice--water coexistence at \(0\,^\circ\mathrm{C}\), boiling water at \(100\,^\circ\mathrm{C}\), and divides the interval into a hundred steps. Kelvin places zero at the limit beyond which no more energy can be extracted, absolute zero, using Celsius-sized steps: \(T = t + 273.15\). Why temperature has a lower end but no upper end returns in Chapter 30.

Temperature brings thermometers. Each selects a property varying monotonically with temperature as its pointer. Mercury and alcohol expand and contract against a glass-tube scale; electronic clinical thermometers use a semiconductor's temperature-dependent resistance; industrial thermocouples use tiny voltages at junctions of dissimilar metals. Infrared thermometers are especially interesting: without contact, they infer temperature from emitted thermal radiation, watching your electromagnetic waves in Chapter 25's language. All share one feature: what they directly display is always their \emph{own} sensing element's temperature or state. The zeroth law allows that reading to be attributed to the object.

\textbf{New concept four: microscopic motion.} What underlies temperature? Nineteenth-century kinetic theory supplied a picture: all matter consists of tiny particles in ceaseless random motion, and temperature measures its \emph{intensity}. Higher temperature means more vigorous microscopic motion. Chapter 27 quantitatively links this to average molecular kinetic energy. Here the qualitative version suffices: heat is not stored fluid, but the macroscopic appearance of particles' ongoing motion.

\textbf{New concept five: internal energy, heat, and work each have one job.} This is the chapter's most important separation.
\begin{itemize}[leftmargin=1.8em]
\item \textbf{Internal energy}: the sum of microscopic particles' kinetic and interaction potential energies inside a system. It is a \emph{state variable}: water's current internal energy depends only on its present state, not how it arrived there.
\item \textbf{Heat}: energy transferred between systems because of a temperature difference. It is a \emph{process quantity}, existing only during transfer. Saying this heating transferred five thousand joules of heat is correct; saying this water contains five thousand joules of heat is a category error, like asking how much raining a reservoir contains. Rainfall is a process; stored water is a state.
\item \textbf{Work}: another route to internal-energy change, driven not by temperature differences but by mechanical pushing, compression, and friction. Rubbing hands and compressing air with a bicycle pump increase internal energy through work, without anyone transferring heat to them.
\end{itemize}
A convenient analogy: internal energy is a bank balance, heat a transfer from someone else, and work overtime wages. The similarity is a distinction between state, the balance, and two transaction routes. The limit is interesting: once money arrives it mixes, just as internal energy retains no label identifying heat or work, so even that part works. But money does not spontaneously flow from rich accounts to poor ones, whereas temperature differences alone determine heat-flow direction. There the analogy fails.
\end{mainline}

\step{The Model in One Sentence}

\begin{oneline}
Temperature labels a state: equilibrated objects share it, and it measures random microscopic motion. Internal energy is the microscopic energy account, a state variable. Heat is energy crossing a boundary because of a temperature difference, a process quantity. Heat can be transferred, produced, or transformed, but never contained. Saying an object contains heat is like saying an account contains bank transfers.
\end{oneline}

Check each phenomenon. The middle cup has one state and one temperature, yet opposite hand reports, so hands do not report temperature. Equal-temperature metal and wood feel different, so the hand reports the \emph{rate of energy loss}. A spark has high temperature but tiny mass and internal energy, hence little transferable heat. Rubbing hands adds work to internal energy with no heat involved. One sentence covers them all.

\step{The Mathematical Translator}

With intuition established, translate it into calculation. The central question: how much does temperature change when a given amount of heat is transferred?

\textbf{First formula: specific heat capacity.} Experiments show that without a phase change, absorbed heat \(Q\) is proportional to mass \(m\) and temperature change \(\Delta T\):
\[
Q \;=\; c\, m\, \Delta T,
\]
The coefficient \(c\), specific heat capacity, is a material property: heat required to raise one kilogram by one kelvin. Water's famously large value, \(c_{\text{water}} \approx 4.2\times10^{3}\ \mathrm{J/(kg\cdot K)}\), is nine times iron's roughly \(0.45\times10^{3}\). Equal heat barely raises generous-capacity water's temperature but warms iron noticeably. Smaller coastal than desert day--night temperature swings, water-cooled engines, and traditional hot-water bottles all reflect that number.

Check special cases immediately. Set \(\Delta T = 0\): no temperature change means no heat absorbed or released, reasonable without phase change. Make \(c\) very large: the same \(Q\) yields tiny warming, explaining water rather than hot iron in a hot-water bottle. Make \(Q\) negative: heat leaves and temperature falls, with signs automatically handling direction.

\textbf{Second formula: the equilibrium account.} Mix hot and cold water in an insulated cup. Heat released by the hot water equals heat absorbed by the cold, caloric theory's accidentally correct domain: without work or phase change, heat behaves as a conserved account. Let hot-water mass be \(m_1\), initial temperature \(t_1\), cold-water mass \(m_2\), initial temperature \(t_2\), and final temperature \(t\):
\[
c\,m_1\,(t_1 - t) \;=\; c\,m_2\,(t - t_2).
\]
One kilogram at eighty degrees mixed with one kilogram at twenty gives \(t = 50\,^\circ\mathrm{C}\), exactly the average. Check: make \(m_2\) much larger than \(m_1\), and \(t\) nearly equals \(t_2\). A boiling drop scarcely changes a swimming pool, reasonably.

\textbf{Third formula: latent heat, where the thermometer stops.} Heat an ice--water mixture continuously. The thermometer remains fixed at zero until the last ice melts, then starts rising. Heat enters but temperature stays still, impossible according to \(Q=cm\Delta T\), showing its no-phase-change condition has failed. During melting, heat dismantles ice's microscopic structure rather than intensifying motion. This is \emph{latent heat}:
\[
Q \;=\; mL,
\]
where \(L\) is latent heat per unit mass. Ice's fusion value is about \(3.34\times10^{5}\ \mathrm{J/kg}\); water's vaporization value at a hundred degrees is about \(2.26\times10^{6}\ \mathrm{J/kg}\). Notice the units: no \(\Delta T\), because temperature stays constant during the phase change. This explains why steam at a hundred degrees burns much worse than boiling water at the same temperature. Condensing on skin, each gram first releases over two thousand extra joules of latent heat, then continues releasing heat as hot water. Conversely, evaporation must \emph{absorb} latent heat, explaining cooling sweat and shivering in a breeze after leaving a pool: faster evaporation pumps heat from your skin.

\begin{center}
\begin{tikzpicture}[scale=1.0]
  % Heating curve: ice, water, vapor plateaus
  \draw[-{Stealth},thick] (0,0) -- (10.2,0) node[below left=2pt and 0pt] {};
  \node[below] at (9.6,-0.05) {\small Heat absorbed \(Q\)};
  \draw[-{Stealth},thick] (0,0) -- (0,4.6) node[above] {\small Temperature};
  % Ticks
  \draw (0,1.15) -- (-0.12,1.15) node[left] {\small \(0\,^\circ\mathrm{C}\)};
  \draw (0,3.35) -- (-0.12,3.35) node[left] {\small \(100\,^\circ\mathrm{C}\)};
  % Curve
  \draw[very thick,blue!60!black] (0,0.25) -- (1.4,1.15);
  \draw[very thick,blue!60!black] (1.4,1.15) -- (3.2,1.15);
  \draw[very thick,blue!60!black] (3.2,1.15) -- (5.2,3.35);
  \draw[very thick,blue!60!black] (5.2,3.35) -- (8.6,3.35);
  \draw[very thick,blue!60!black,dashed] (8.6,3.35) -- (9.7,4.3);
  % Labels
  \node[below] at (0.7,0.35) {\small Ice warms};
  \node[above] at (2.3,1.15) {\small Ice melts (constant temperature)};
  \node[above,rotate=32] at (4.2,2.1) {\small Water warms};
  \node[above] at (6.9,3.35) {\small Water vaporizes (constant temperature)};
  \node[above,rotate=38] at (9.15,3.95) {\small Steam warms};
\end{tikzpicture}
\end{center}

The two plateaus represent latent heat: heat enters while temperature stays still. Vaporization's plateau is much longer than melting's. Vaporizing a kilogram of water at a hundred degrees takes over five times the heat needed to bring a kilogram from zero to boiling.

\textbf{An estimate with real numbers.} How long does a \(1500\ \mathrm{W}\) kettle need to boil a liter, one kilogram, of water initially at twenty degrees Celsius?
\[
Q = c\,m\,\Delta T = 4.2\times10^{3}\times 1\times (100-20) \approx 3.4\times10^{5}\ \mathrm{J}.
\]
Power \(1500\ \mathrm{W}\) supplies fifteen hundred joules per second, giving an ideal time of about \(3.4\times10^{5}/1.5\times10^{3}\approx 220\) seconds, three or four minutes, matching kitchen experience. A real kettle takes longer because its body and the air take some heat. Reframe the same number: 340,000 joules could melt a whole kilogram of ice at zero into water at zero, or raise ten kilograms of water by eighty degrees. Energy accounts never equivocate about numbers.

\textbf{Fourth: heat's three routes.} Heat moves from higher to lower temperature in only three ways.
\begin{itemize}[leftmargin=1.8em]
\item \textbf{Conduction}: microscopic collisions relay vigorous motion from particle to particle without the particles relocating. Metals relay efficiently; wood, air, and down slowly. Roughly, greater temperature difference, larger contact area, and shorter distance mean faster transfer.
\item \textbf{Convection}: heating changes fluid density and drives bulk flow, carrying energy along with the particles themselves. A radiator warms a room chiefly through the air circulation it drives, not its modest radiation.
\item \textbf{Thermal radiation}: objects at any temperature emit electromagnetic energy, more strongly when hotter. No medium is needed, so vacuum offers no barrier. This is how the Sun warms Earth, the same physics as Chapter 25's waves.
\end{itemize}
A vacuum flask engineers barriers to all three. Vacuum between double walls blocks conduction and convection, since no medium means neither relays nor relocation. Silvered inner walls reflect radiation, and a cork or plastic lid closes the remaining conduction leak. Down clothing chiefly traps abundant still air in fluffy clusters. Still air conducts poorly; down does not generate warmth but minimizes your heat-loss rate. Two more convection examples: coastal daytime wind nearly always blows inland because land's lower heat capacity lets it warm quickly, with rising hot air replaced by cooler sea air. This circulation is fluid relocating. A convection oven's fan likewise cooks food more evenly. Radiation offers a counterintuitive example: beside a large window on a deep-winter night, even a warm room can feel as if the window sucks heat away. It is not cold air entering, but your skin radiating through the glass toward the cold outdoors, losing energy at light speed.

\begin{mathbridge}
Radiation's quantitative rule is strikingly elegant. A surface emits power
\[
P \;=\; \varepsilon \sigma A T^{4},
\]
where \(T\) must be in kelvins, \(A\) is surface area, \(\sigma \approx 5.67\times10^{-8}\ \mathrm{W/(m^{2}\cdot K^{4})}\) is the Stefan--Boltzmann constant, and \(\varepsilon\) is emissivity between zero and one, near one for dark rough surfaces and near zero for shiny polished silver. The fourth power is fierce: double temperature and radiation grows sixteenfold. This explains why you cannot approach a steel furnace and why a small change in Earth's mean temperature greatly changes its outgoing energy account. Conduction is similarly compact: heat per unit time through material is proportional to temperature difference and area, inversely proportional to thickness, with coefficient called thermal conductivity. Metals can conduct thousands of times better than air. Incidentally, any object's emitted spectrum depends only on temperature. This blackbody-radiation rule resisted late-nineteenth-century theory, and the crisis forced the birth of quantum theory. We return to that accident scene later.
\end{mathbridge}

\step{The Model's Limits}

Every model must define its jurisdiction. Five boundaries matter here.

\textbf{First: temperature belongs to large populations.} Its measure of microscopic-motion intensity is an \emph{average}. Across the vast number of particles in a cup of water, it is rock-steady; for a system of a dozen particles, the average fluctuates wildly and temperature has only a fuzzy answer. Asking one molecule's temperature is meaningless, like asking the average height of one spectator. Temperature and pressure are collective concepts.

\textbf{Second: a thermometer reports its own temperature.} A clinical thermometer stays under your arm for minutes not because it reacts slowly, but because the zeroth law needs time to take effect. It must equilibrate with your body before its reading equals your temperature. Infrared thermometers follow another logic: without contact, they infer temperature from skin radiation. Their reading is a radiation-based estimate, so glass barriers and reflective metal surfaces cause errors.

\textbf{Third: temperature may vary point by point.} Ten and twenty centimeters above a flame can differ by hundreds of degrees; a sunlit car roof and its underside also differ. Temperature must be defined locally, with a single shared value only inside an equilibrated region. A satellite's sunlit and shaded sides may differ by over a hundred degrees, so its temperature requires specifying where.

\textbf{Fourth: specific heat is not constant.} In \(Q=cm\Delta T\), \(c\) is approximately constant only over a limited range. At very low temperatures, materials' heat capacities fall sharply toward zero for purely quantum reasons. Together with the fifth boundary, this reminds us that our heat-capacity model approximates ordinary temperatures, not the ultimate law.

\textbf{Fifth: absolute zero may be approached, never reached.} Kelvin zero is no ordinary cold point, but a lower limit to energy extractable from microscopic motion. Theory and experiment agree that every cooling method can approach it step by step but never attain it. The complete reason requires thermodynamics' third law, explained in Chapter 30.

\begin{challenge}
The modern definition goes deeper than degree of hotness. Statistical physics defines temperature as the reciprocal of internal energy's rate of change with entropy. This has a startling legitimate consequence: in special systems with an upper energy bound, such as nuclear spins in an external magnetic field, pumping more particles into high-energy than low-energy states can produce \emph{negative} temperature. Moreover, negative temperatures are hotter than every positive temperature: contact sends heat from negative to positive. Population inversion, the working basis of lasers, is such a state. The example reminds us that an everyday hot--cold scale is the concept's beginning, not its endpoint.
\end{challenge}

\step{Explaining the Opening Phenomenon}

Test the new model by settling the three opening mysteries.

\textbf{Why does iron feel colder than wood?} Their temperatures match. The thermometer does not lie; the zeroth law ensures the room equilibrates overnight. The false assumption is hand equals thermometer. Specialized skin nerves report not current skin temperature but the \emph{rate at which heat leaves it}. Iron's high conductivity quickly removes heat at contact and spreads it through the leg, abruptly cooling local skin and making nerves shout cold. Wood conducts poorly; a little transferred heat remains near its surface, skin cools slowly, and it feels mild. Same temperature, different sensation: transfer rate makes the difference. The hand is a crude \emph{heat-flow meter}, not a thermometer. That is no embarrassment; Chapter 28 shows even precise calorimetry began by tracking heat flow.

% BEGIN TRANSLATOR SAFETY NOTE
\textbf{Translator safety note:} This explanation repeats the source's unsafe hand-soaking premise; do not reproduce it with hot water. Use thermometer readings instead. See \href{https://www.cpsc.gov/s3fs-public/5098.pdf}{CPSC scald guidance}.
% END TRANSLATOR SAFETY NOTE

\textbf{Why do three cups contradict themselves?} The same reason. A hand in hot water for a minute warms until outward heat loss balances. Put it in warm water, cooler than its skin, and heat flows outward: the meter reports cool. The other hand emerges colder than the middle water, so heat enters and it reports hot. The two meters have been calibrated into opposite states, naturally disagreeing on one cup. The zeroth law supplies arbitration independent of touch: thermometers need no hand-soaking ritual.

% BEGIN TRANSLATOR SAFETY NOTE
\textbf{Translator safety note:} Small sparks are not reliably harmless. Keep them off skin and eyes; do not use the source's energy argument to omit protective equipment. See \href{https://www.cpsc.gov/node/65701}{CPSC sparkler guidance} and \href{https://www.osha.gov/etools/shipyard/general-requirements/ppe/ppe-selection/hot-work}{OSHA hot-work guidance}.
% END TRANSLATOR SAFETY NOTE

\textbf{Why do sparks not injure us?} This beautifully separates temperature, internal energy, and heat. A spark really can exceed a thousand degrees, with vigorous microscopic motion. Yet each is a tiny metal or oxide particle, perhaps a ten-thousandth of a gram, with minuscule internal energy. On skin it transfers barely enough heat to warm a pinprick. A kettle is only a hundred degrees, but a kilogram of water has a huge internal-energy account and can deliver tens of thousands of joules to skin. Injury depends on transferred \emph{heat}, not spark \emph{temperature}. That is the title's meaning: temperature is not heat.

For the initial guesses: first, equal temperature, since sensation and temperature differ. Second, the question is a category error; ask instead which has more internal energy, the kettle by far, and which can transfer more heat to skin, again the kettle. Third, while ice remains, phase change pins the mixture at zero. Incoming room heat goes entirely into latent heat without moving the temperature.

Settle the three later observations too. A clinical thermometer needs minutes because equilibration is not instantaneous: it first reaches your temperature while reporting its own. Rice gets hot while its bowl stays warm because microwaves are mainly absorbed by water molecules, as Chapter 25 discussed. Rice contains much water and ceramic little, so their absorbed energies differ greatly; some bowl warming is borrowed by conduction from the rice. Foggy breath illustrates latent heat: water vapor in exhaled air condenses on cold glass into tiny liquid droplets. The white mist is liquid beads, not gas. Seconds later they evaporate, drawing vaporization heat from the glass, and the mist disappears.

\step{Thought Experiments and Challenges}

\begin{thought}[Throw a Thermometer into Deep Space]
Take a sufficiently sturdy thermometer beyond the solar system, far from any star, and wait long enough. What will it read? Without air, neither conduction nor convection exists, but radiation always remains. It radiates energy and cools while receiving radiation from every direction. Eventually it equilibrates with the universe's radiation field at about \(2.7\ \mathrm{K}\), the temperature of the microwave background left by the Big Bang. This shows two things. Radiation allows thermal equilibrium across vacuum, so the zeroth law needs no contact. And the temperature ruler is fantastically long: laboratory cold atoms have reached below a billionth of a kelvin; massive stellar cores exceed ten million kelvins; momentary accelerator-collision temperatures exceed a trillion kelvins. From lowest to highest, over twenty orders of magnitude host entirely different forms of matter.
\end{thought}

\begin{puzzles}
\item Equal masses of iron at eighty and twenty degrees equilibrate at fifty. Replace the cooler iron by an equal mass of \emph{water} at twenty. Is equilibrium above or below fifty? (Hint: compare heat capacities and ask which resists temperature change more. Far below fifty, in fact around thirty degrees.)
\item On the Tibetan Plateau, water boils readily but noodles remain undercooked. Is the flame too weak? Would more flame help? (Hint: boiling is a phase-change plateau, lower at reduced pressure. Stronger heating cannot raise water above its boiling point; extra heat leaves with steam. A pressure cooker deliberately raises that plateau.)
\item After an hour in summer sunlight, will iron and wooden railings feel more different or less different than in winter? (Hint: analyze your hand as a heat-flow meter, with flow reversed. Are sunlit wood and metal actually at equal temperature? Radiation and conduction together make the answer more complex than in winter.)
\item Put ice at zero degrees into water at zero and close an insulating lid. Will it melt? What if the whole cup sits in a room at zero? (Hint: first check for a temperature difference. None means no heat flow. Then ask where latent heat would come from.)
\end{puzzles}

We have separated temperature, heat, and internal energy, laying out the ledger without its governing rules. How do heat and work together determine internal-energy change? Why are some accounts possible and others forever impossible? Chapters 28 and 29 establish those two laws.
